\chapter{Methodology}
\label{sec:methodology}
The primary objective of this research is to port and extend AnitoPlume's real-time simulation framework by modernizing its rendering pipeline, optimizing hardware utilization through new, explicit graphics APIs, and refining its underlying physical simulation models. 

\section{Implementation and Migration}
\subsection{Legacy Simulation}
The legacy execution architecture relies on a strictly sequential, single-threaded pipeline driven by an iterative time-stepping loop. 
Within each frame slice, the simulation runs a high-frequency sub-stepping loop (executing ten physics updates per render frame) to remove visual smoke jumping and oscillations due to numerical instability. 
The control flow inside this sub-step operates linearly across distinct global entity arrays. 

As shown in Figure \ref{fig:legacy-simulation}, the simulation is organized around the smoke layer as the primary simulation unit. Each smoke layer models the behavior of the eruption column by computing atmospheric air entrainment, thermodynamic properties (e.g., temperature and density), external forces such as buoyancy, gravity, wind, and drag, and integrating the resulting motion over time. Attached to every smoke layer are free spheres, which represent smoke that inherit the motion and physical state of their parent layer while evaluating transitions into falling or neutrally buoyant states. Each free sphere further generates a hierarchy of subspheres that provide the volumetric detail required for visualization. These subspheres are rendered as billboard sprites to produce the final volumetric appearance of the smoke cloud but do not participate in the physical simulation.

\begin{figure}[H]
    \centering
    \makebox[\textwidth][c]{
        \includegraphics[width=0.9\textwidth]{legacy-sims2.png}
    }
    \caption{Legacy physics simulation loop.}
    \label{fig:legacy-simulation}
\end{figure}

AnitoPlume initializes the simulation environment by loading a preprocessed three-dimensional terrain model representing Taal. The terrain mesh is generated offline from geospatial elevation datasets and optimized for interactive rendering while preserving the topographic features required for volcanic hazard visualization \cite{del2026interactive}. 
Textures for the terrain model were made through photobashing of satellite images taken from Google Earth; as seen on Figure \ref{fig:photobashing} \cite{del2026interactive}. 

\begin{figure}[H]
    \centering
    \makebox[\textwidth][c]{
        \includegraphics[width=0.8\textwidth]{photobash.png}
    }
    \caption{Workflow for texturing the surrounding area of Taal Volcano. From \protect\cite{del2026interactive}}
    \label{fig:photobashing}
\end{figure}

During initialization, the terrain is imported into the rendering pipeline, transformed into the simulation coordinate system, and associated with its corresponding surface texture.

\begin{figure}[H]
    \centering
    \makebox[\textwidth][c]{
        \includegraphics[width=0.9\textwidth]{textureRender.png}
    }
    \caption{Legacy terrain rendering pipeline.}
    \label{fig:texture-rendering}
\end{figure}

In addition to visualization, the terrain serves as the physical boundary of the simulation. Height and surface normal information are queried during particle updates to detect terrain intersections and compute particle-ground interactions, allowing falling volcanic material to respond to the underlying topography. Consequently, the terrain functions as both the rendered environment and the collision surface used by the simulation.

Beyond the core simulation and terrain rendering pipelines, AnitoPlume also provides an interactive graphical user interface (GUI) through which users can manipulate plume behavior in real time. This interface additionally surfaces contextual 3D tooltips that display relevant information about Taal Volcano directly within the scene. Figure~\ref{fig:uirender} shows this interface, representing the primary means by which users interact with and observe the legacy simulation system.

\begin{figure}[H]
    \centering
    \makebox[\textwidth][c]{
        \includegraphics[width=0.95\textwidth]{UI3.png}
    }
    \caption{GUI of AnitoPlume for manipulating plume behavior and 3D UI tooltips for displaying Taal Volcano information to the user. From \protect\cite{del2026interactive}}
    \label{fig:uirender}
\end{figure}

\subsection{Initial AnitoPlumeV2 Prototype}

Prior to this study, an early prototype of AnitoPlume V2 was developed as an exploratory effort to modernize the visual rendering of the Taal Volcano model. However, the prototype was discontinued before the full implementation of plume simulation, terrain rendering, and user interaction, could be completed. Because the prototype represents the starting point from which this work proceeds, its implemented capabilities and architectural limitations must be clearly established.

In its current initial state, the prototype renders a single textured terrain mesh of Taal Volcano under a ray traced lighting pipeline, navigable through a first person camera interface. No plume simulation, volumetric ash visualization, particle system, interactive GUI, or terrain collision logic has been integrated at this stage. The prototype functions exclusively as a static scene viewer for the volcano model and serves as the starting point that this study will extend into a full simulation and visualization system.

The V2 prototype is built upon NVIDIA Falcor \cite{falcor2024}, a real-time rendering research framework that provides a structured pipeline for high-fidelity scene visualization. The current implementation focuses exclusively on the visual rendering of the volcano surface; no physics-based plume simulation, volumetric ash dispersion, or particle-based eruption dynamics have been implemented at this stage. 

\begin{comment}
The rendering pipeline is composed of several passes executed in sequence:

\begin{itemize}
    \item \textbf{Visibility Buffer Ray Tracing (VBufferRT).} This pass leverages shader model 6.5 and raytracing tier 1.1 to generate a packed visibility buffer. For each pixel, the pass stores primitive indices and barycentric coordinates, along with auxiliary data such as depth values, motion vectors, and view-space information. The visibility buffer serves as the geometric foundation for subsequent lighting and accumulation passes.
    
    \item \textbf{RTXDI Direct Illumination.} The visibility data is processed using NVIDIA's RTXDI (Ray Tracing eXtensible Direct Illumination) algorithm, which computes diffuse and specular lighting contributions on all visible surfaces. This pass replaces the simple diffuse lighting used in the legacy AnitoPlume V1 with a physically motivated direct illumination model.
    
    \item \textbf{Temporal Accumulation (AccumulatePass).} To reduce rendering noise introduced by the stochastic ray tracing process, the AccumulatePass performs temporal frame accumulation. Configurable precision modes allow the integration of per-pixel color values across multiple frames, progressively converging toward a noise-free result.
    
    \item \textbf{Tone Mapping.} The final pass converts the accumulated high dynamic range (HDR) color output into displayable low dynamic range (LDR) values, applying configurable parameters for exposure compensation, automatic exposure correction, white balance adjustment, and tone mapping operators.
\end{itemize}
\end{comment}

The scene is navigated through a first-person camera interface with configurable movement speed and supports standard scene hierarchy management, including environment map assignment, light source definition, and material property authoring. Figure~\ref{fig:v2-prototype} illustrates the visual output of the prototype, showing the Taal Volcano terrain rendered with the ray tracing and temporal accumulation pipeline described above.

\begin{figure}[H]
    \centering
    \makebox[\textwidth][c]{
        \includegraphics[width=0.95\textwidth]{v2-prototype.png}
    }
    \caption{Current initial state of the AnitoPlume V2 prototype, showing the Taal Volcano terrain rendered through the VBufferRT and RTXDI passes within the Falcor framework. No plume simulation is present at this stage.}
    \label{fig:v2-prototype}
\end{figure}


Despite the visual capabilities of the existing render pipeline, the V2 prototype carries a significant architectural dependency to the NVIDIA ecosystem. 
The study adopts a framework-agnostic approach, meaning the final architecture must not rely on any single rendering framework such as Falcor. 

Consequently, the V2 prototype as provided cannot serve as a guaranteed foundation for the final system. The study treats the existing prototype code as a reference implementation of the visibility buffer pipeline and the accumulated RTXDI logic, but reserves the ability to re-implement the full application from scratch using a framework- and GPU-agnostic explicit graphics API if a more suitable rendering architecture is warranted. This design decision ensures that the methodology outlined in this chapter, namely parallel simulation, explicit API resource management, and cross-vendor benchmarking, is not constrained by the dependencies of the discontinued prototype.

\section{Benchmarking Procedure}
\subsection{Profiling Tool Selection}
\label{sec:profiling_tool_selection}

As discussed in Section~\ref{sec:profiling_tools}, profiling tools for explicit graphics API development fall into several categories: real-time CPU-GPU profilers, hardware-level performance analyzers, and frame capture debuggers. Each offers a different scope of instrumentation.
The selection of a primary profiling instrument for this study requires that the tool satisfy three conditions simultaneously:

\begin{enumerate}
    \item \textbf{Cross-vendor hardware support.} The benchmarking procedure is designed to produce reproducible results across a range of consumer and workstation GPUs from NVIDIA, AMD, and Intel. A profiling instrument locked to a single GPU vendor would limit the generalizability of the results and introduce hardware-specific reporting biases.
    \item \textbf{Framework-agnostic comparison.} Because this study evaluates both the legacy OpenGL implementation in Phase~1 and the ported explicit API pipeline in Phase~3, the benchmarking methodology must remain independent of any specific rendering framework. Both pipelines will be profiled under identical hardware and workload conditions, and the same profiling tool must capture both runs using the same metric definitions and output formats. This approach ensures that performance comparisons reflect architectural differences rather than artifact of tooling or framework-specific optimizations.
    \item \textbf{Evaluation metrics.} The tool must expose per-event GPU execution time, shader-stage latency, and memory bandwidth usage to populate the metrics defined in Table~\ref{tab:performance_metrics}.
\end{enumerate}

\textbf{RenderDoc} \cite{renderdoc2023} is selected as the primary profiling instrument for this study. 
While Tracy (Section~\ref{sec:tracy}) identifies multi-frame trends such as thread synchronization and CPU-GPU overlap, it records only aggregated zone-level timing and cannot isolate per-event draw calls or inspect individual pipeline states. RenderDoc's offline frame-capture model provides the per-event granularity that Tracy lacks which is necessary to populate the metrics in Table~\ref{tab:performance_metrics}, particularly Shader Execution Time, Command Recording Time, and GPU Memory Bandwidth. 
RenderDoc also supports the API used by the original AnitoPlume (OpenGL) and modern graphics APIs such as but not limited to Vulkan and DirectX~12 under an identical capture workflow, ensuring that profiles from Phases~1 and~3 are methodologically compatible and directly comparable, and it operates across NVIDIA, AMD, and Intel GPUs, unlike NVIDIA Nsight Graphics (Section~\ref{sec:nsight}) which is restricted to a single vendor and would limit the generalizability of benchmarking results.
\subsection{Evaluation Metrics}
To evaluate scalability, performance metrics will be collected through increasing simulation workloads, using profiling tools such as RenderDoc to capture low-level pipeline data alongside application-level timing instrumentation. 
The specific metrics utilized for this evaluation, along with their respective descriptions, are detailed in Table~\ref{tab:performance_metrics}.

\begin{table}[htbp]
\centering
\renewcommand{\arraystretch}{1} % Set to 1.3 for a clean, professional breathing room
\caption{Performance metrics and descriptions used for evaluation.}
\label{tab:performance_metrics}
\begin{tabular}{l p{8cm}}
\hline
\textbf{Metric} & \textbf{Description} \\ 
\hline
Average FPS            & Overall rendering performance \\
Frame Time (ms)        & Total frame execution time \\
CPU Frame Time         & CPU workload \\
GPU Frame Time         & GPU workload \\
Command Recording Time & CPU time spent preparing GPU commands \\
Shader Execution Time  & GPU shader workload \\
GPU Memory Bandwidth   & Memory utilization \\
CPU Utilization        & Effectiveness of multithreading \\
\hline
\end{tabular}
\end{table}

Workload size and rendering performance will be visualized using line graphs illustrating trends in average frame rate, frame time, CPU utilization, and GPU utilization. GPU profiling metrics, including shader execution time, command recording time, and memory bandwidth utilization, will likewise be compared to determine how efficiently the proposed architecture utilizes available hardware resources. 

Finally, the benchmarking results will be interpreted alongside the User Acceptance Testing (UAT) results to determine whether improvements in computational performance correspond to improvements in perceived responsiveness and overall user experience.


\section{Usability Testing}

To verify that the architectural overhaul maintains or improves the interactive user experience, an exploratory usability test will be conducted with target beneficiaries:
\begin{itemize}
    \item Participants: Recruitment will follow a tiered priority order. PHIVOLCS personnel will be the top priority given their domain expertise in volcanology and direct familiarity with Taal Volcano hazard assessment. Should this pool be insufficient, recruitment will expand to, in order: (1) individuals who have witnessed a Taal Volcano eruption firsthand, (2) enthusiasts with experience using volcano/geohazard simulation software, and (3) computer science students with a background in rendering or volcano-related coursework/projects. This staged recruitment reflects the tool's intended end-users while ensuring an adequate sample size, and informs the corresponding participant categories documented in the study's ethics clearance and consent materials.
    \item Setup:  After screening and obtaining consent from the Informed Consent Form (ICF), participants will be given time to explore the interface.
    \item Task:  Participants will use the software to recreate specific Taal Volcano eruption plumes, following a set of predefined scenarios.
    \item Survey: Users will complete the System Usability Scale (SUS) questionnaire to rate the software's ease of use.
    \item Feedback: A closing interview will gather suggestions for improving the interface and overall performance.
\end{itemize}

% Changed to [H] to force exact placement before the next section
\begin{table}[H]
\centering
\caption[Sample SUS questionnaire.]{Sample SUS questionnaire implemented by \cite{del2026interactive}. 1 --- strongly disagree, 5 --- strongly agree.}\begin{tabular}{lp{12cm}}
\hline
Code & Question \\
\hline
Q1 & I think that I would use AnitoPlume frequently. \\
Q2 & I found AnitoPlume unnecessarily complex. \\
Q3 & I found AnitoPlume easy to use. \\
Q4 & I think that I would need the support of a technical person to be able to effectively and efficiently use AnitoPlume. \\
Q5 & I found that various features of the simulator were well integrated in AnitoPlume. \\
Q6 & I thought there was too much inconsistency in AnitoPlume. \\
Q7 & I would imagine that most people without prior experience in using similar software would easily learn how to use AnitoPlume. \\
Q8 & I found AnitoPlume very cumbersome (slow or complicated and therefore inefficient) to use. \\
Q9 & I felt very confident while using AnitoPlume. \\
Q10 & I needed to figure out a lot of things (such as hotkeys, navigating the UI) before I could confidently use AnitoPlume. \\
\hline
\end{tabular}
\end{table}


\section{Research Workflow}

This study follows a three-phase research workflow, outlined in Table \ref{tab:research_phases}, that guides the porting and extension of AnitoPlume from its legacy rendering implementation to the proposed rendering architecture, from initial legacy analysis through final evaluation of performance and usability.


\begin{table}[htbp]
\centering
\small
\renewcommand{\arraystretch}{1.2}
\begin{tabular}{l p{3.5cm} p{6.5cm}}
\hline
\textbf{Phase} & \textbf{Research Activity} & \textbf{Purpose / Output} \\
\hline
\textbf{1} & \textbf{Legacy Pipeline Profiling and Bottleneck Analysis} & Evaluate the existing simulation and rendering architecture to establish baseline performance metrics. Profile CPU-GPU synchronization, draw call overhead, and frame pacing to identify compute and graphics pipeline bottlenecks that impede scalability. \\
\hline
\textbf{2} & \textbf{Architectural Refactoring and Pipeline Optimization} & Design and iteratively implement optimizations within the execution model and rendering pipeline. This involves exploring modern graphics paradigms (e.g., explicit command submission, reducing driver overhead) and evaluating concurrent execution strategies to maximize render throughput, adapting the architecture based on empirical performance gains. \\
\hline
\textbf{3} & \textbf{Performance Benchmarking and Interactive Evaluation} & Conduct a comparative analysis between the legacy and refactored pipelines using low-level GPU profiling metrics (e.g., shader execution time, memory bandwidth utilization). Perform user testing to correlate these hardware-level optimizations with perceived interactive responsiveness and visual consistency. \\
\hline
\end{tabular}
\caption{Overview of the research phases, activities, and objectives.}
\label{tab:research_phases}
\end{table}


\begin{comment}

\begin{table}[]
    \centering
    \begin{tabular}{l|c|c|c|c|c|c|c|c|c}
         \toprule
         & \multicolumn{6}{c|}{2021} & \multicolumn{3}{c}{2022} \\
         \textbf{Activity} & \textbf{Apr} & \textbf{May} & \textbf{Jun} & \textbf{Jul} & \textbf{Aug} & \textbf{Sep} & \textbf{Jan} & \textbf{Feb} & \textbf{Mar} \\
         \midrule
         Data Collection & **** & **** & & & & & & & \\
         \midrule
         Software Design & & **** & **** & **** & **** & **** & & & \\
         \midrule
         Validation & & & & & & **** & **** & **** & \\
         \bottomrule
    \end{tabular}
    \caption{Gantt Chart of Activities}
    \label{tab:timetableactivities}
\end{table}
\end{comment}